% !TEX root = ../main.tex

\section{Related Literature}
\label{sec:litreview}

Standard portfolio theory suggests that risk-averse investors should hold some portion of their wealth in a diversified portfolio of equities to earn the equity premium. However, the actual stock market participation (SMP) rate is significantly lower than predicted, a puzzle widely documented in the foundational literature. Subsequent studies have explored various determinants to explain this participation puzzle. \citet{menkhoff2025} provide a comprehensive overview, categorizing the drivers into direct costs such as wealth and income \citep{guiso2008}, knowledge and cognitive abilities \citep{vanrooij2011}, individual trust \citep{guiso2008, bu2022}, and peer effects through sociability and networks \citep{hong2004, brown2008}. 

In this paper, I focus on a yet unexplained pattern: the persistent geographic variation in SMP across U.S. states, which cannot be fully accounted for by differences in income \citep{chien2017}. Prior literature has used state-level variances to isolate driving factors of SMP. \citet{brown2004} find that individuals in states with higher average SMP are more likely to participate themselves, suggesting a community ownership effect. \citet{giannetti2016} use a 2002 corporate fraud scandal to show that corporate fraud revelation leads to lower SMP and decreased stock holdings in all firms, suggesting a local shock to trust in the stock market. Also suggestive of a trust channel, \citet{huang2023} find that states with higher accounting quality, proxied by audit market size, exhibit higher SMP.  

Another crucial driver of stock market participation is the political orientation of the investor. Using Finnish data, \citet{kaustia2011} find that right-leaning individuals are more likely to participate in the stock market. However, when extended to other European countries, \citet{kaustia2023} find that the relationship between right-leaning political orientation and SMP is not consistent, particularly in countries with low regulatory quality. In the United States, research has shown an alignment effect: individuals become more optimistic and perceive markets as less risky when their preferred political party is in power, which consequently boosts their likelihood to invest \citep{bonaparte2017}. Building on this, \citet{meeuwis2022} demonstrate that likely-Republicans increased the equity share of their portfolios relative to Democrats immediately following the 2016 presidential election.  

While \citet{meeuwis2022} focus on the immediate, short-term rebalancing of equity shares within existing brokerage accounts around a specific election event, this paper examines baseline stock market participation rates during normal times. Furthermore, while previous studies have often relied on U.S. data that is fragmented across different time periods for stock market variables versus political indicators, this study employs unified datasets to test the state-level political stance and its dual-channel effect on SMP. 